% ── CHAPTER 5: WHY QUANTUM MECHANICS POINTS TO PROCESS ONTOLOGY (~8,000 words) ──
\chapter{Why Quantum Mechanics Points to Process Ontology}
\label{ch:qm-requires-process}

\epigraph{What is proved by impossibility proofs is lack of imagination.}{John Bell}

\firstword{T}{here} is a temptation, when first encountering quantum mechanics, to treat it as
a collection of paradoxes: the cat that is simultaneously alive and dead, the
particle that goes through both slits at once, the measurement that seems to
disturb what it observes.
This temptation should be resisted.
Quantum mechanics is not primarily a collection of puzzles awaiting a future
resolution.
It is a mature physical theory with a precise mathematical structure and an
experimental record, now a century long, in which none of its predictions has
been contradicted.
Its puzzling features are not failures of the theory.
They are accurate descriptions of how nature actually behaves.

The philosophical challenge is not to explain away the puzzles but to ask what
kind of world makes a theory like this true.
That is the question this chapter pursues.
Its answer has a precise logical structure.

The argument runs in four steps.
First, Bell's theorem, confirmed experimentally from Aspect's 1982 tests through
the 2022 Nobel Prize, eliminates a specific and fundamental class of theories:
those that assign definite local properties to quantum systems before measurement.
Second, if particles have no definite local properties before measurement, the
best account of what they have instead is relational: properties that come into
being through interaction rather than being possessed independently.
Third, substance ontology --- the framework that has governed Western metaphysics
from Aristotle through Newton and beyond --- requires, in the local form that
classical physics gave it, exactly what Bell eliminates.
That form of substance ontology is therefore not merely challenged or complicated
by quantum mechanics; it is ruled out by a mathematical theorem and its
experimental confirmation, given assumptions I will state.
Fourth, process ontology --- the framework Whitehead developed in 1929 and that
has been extended by philosophers of physics ever since --- is the most fully
developed alternative.
The structural alignment between process ontology and quantum mechanics is, I
will argue, not accidental: it is point-by-point, covering the four features of
quantum mechanics that are hardest to explain within any framework.

I want to be careful about the strength of the claim.
The argument is not that process ontology has been proven true.
It is that local substance ontology has been formally ruled out, and that among
the surviving frameworks, process ontology has the strongest systematic claim.
That asymmetry is philosophically significant even in the absence of proof.
The four steps differ in kind.
The first and third are deductive: grant their premises, and their conclusions
follow.
The second and fourth are inferences to the best explanation.
Quantum mechanics requires, in the strict sense, giving up the picture of
separate things each carrying its own definite properties.
It does not by itself select process ontology over every rival that survives;
that is a judgment of best fit, and I will defend it as one.
By the end of this chapter, the reader will have the conceptual tools to understand
why the following conclusion is not an overstatement: quantum mechanics does not
merely \emph{suggest} process ontology.
Given the experimental record, it \emph{requires} something like it: an ontology
in which relations do work that the intrinsic properties of separate things
cannot do.

The epigraph is a warning I mean to heed.
Bell wrote it in defense of the pilot-wave theory, against proofs, von Neumann's
among them~\citep{vonNeumann1932}, that had been taken to show hidden variables
impossible.
An argument from a no-go theorem owes its reader a plain statement of what the
theorem assumes and which rivals it leaves standing.


% ─────────────────────────────────────────────────────────────────────────────
\section{Step 1: Bell's Theorem --- What It Actually Proves}
\label{sec:ch5-step1}
% ─────────────────────────────────────────────────────────────────────────────

To understand what Bell's theorem proves, it helps to begin with what Einstein,
Podolsky, and Rosen thought they had established in 1935.

Their paper~\citep{Einstein1935EPR} made the following argument.
Quantum mechanics predicts that two particles, once they have interacted,
can remain correlated in a way that seems to transcend distance: measuring one
particle immediately determines what the corresponding measurement on the other
would show, even if the two particles are on opposite sides of the galaxy.
Einstein, Podolsky, and Rosen found this deeply objectionable.
No signal can travel faster than light; and yet the correlation exists.
Their proposed explanation was that quantum mechanics is simply incomplete.
The particles carried some additional, hidden information all along --- a
predetermined plan for how they would respond to any measurement --- and the
apparent mystery dissolves when you accept that the theory is just missing those
hidden variables.

This was not an unreasonable position.
It was, in fact, the commonsense scientific position: when a theory produces
surprising correlations, look for an underlying mechanism, some prior agreement
encoded in the particles themselves.

In 1964, John Stewart Bell asked a question that had not been asked before:
could this proposal --- that particles carry hidden local properties that
predetermine their measurement outcomes --- actually reproduce the statistical
predictions of quantum mechanics?
The answer, he showed, was no~\citep{Bell1964}.

Bell's theorem is a mathematical result.
Its two central premises are these.
The first is \emph{locality}: the idea that what happens to one particle cannot
depend on what happens to a distant, spacelike-separated particle.
More precisely, the value that a measurement reveals for particle A cannot depend
on which measurement setting was chosen for particle B, if B is so far away that
no signal could have traveled between them in time to influence the outcome.
This is not a radical assumption.
Newton's gravity violated it, to Newton's own discomfort, but field theory from
Maxwell onward restored it, and special relativity made it a principle.
The second premise is \emph{realism} in a specific, technical sense: the idea that
the particles have definite values for the measured properties \emph{before} the
measurements are performed.
Not that we know those values, but that they exist.
That the particle has, say, a definite spin orientation from the moment it is
created, waiting to be revealed when we look.

Bell proved that any theory satisfying both premises --- locality and prior definite
values --- must obey certain inequalities.
These inequalities place upper bounds on the statistical correlations that can
exist between measurements on distant particles.
They are the CHSH inequalities, named for \citet{CHSH1969}, who put Bell's
original argument into the form most directly testable by experiment.

Quantum mechanics predicts that these inequalities are violated.
The correlations between entangled particles, according to the quantum formalism,
are stronger than anything a locally realistic theory can produce.
Experiment has consistently confirmed the quantum prediction.

\citet{Aspect1982} performed the first decisive experimental test.
Working in Paris with pairs of entangled photons, Aspect, Grangier, and Roger
found clear, unambiguous violations of the Bell inequalities, precisely as
quantum mechanics predicted.
Later the same year Aspect, Jean Dalibard, and Roger switched the measurement
settings while the photons were in flight, a design intended to close the
loophole of any signal traveling between measurement stations.

For two decades after Aspect, skeptics noted that various technical loopholes
remained: perhaps the detector efficiency was low enough that the sample of detected
particles was not representative; perhaps the random settings were not truly
independent of the hidden variables.
Between 2015 and 2017, a series of loophole-free experiments simultaneously
closed both the locality loophole and the detection efficiency loophole for the
first time.
The 2022 Nobel Prize in Physics was awarded jointly to Alain Aspect, John F.
Clauser, and Anton Zeilinger for this body of experimental work~\citep{Zeilinger2023Nobel}.

What does the experimental record actually prove?

It proves that no theory satisfying both locality and prior definite values can
reproduce the predictions of quantum mechanics, given two further assumptions
set out at the end of this section.
In other words: either quantum particles do not have definite values before
measurement, or those values are not local.
One of the two premises must be abandoned.

Many physicists describe the result as the failure of \emph{local realism} and
take the first horn, and so do I.
Quantum particles do not carry definite values for their observables prior to
measurement.
The measurement is not the passive reading of a pre-existing fact.
Something genuinely new comes into being through the act of measurement.
The theorem alone does not force this choice, though, and Bell himself did not
make it.

I want to dwell on the word ``local'' because it carries more philosophical weight
than is often appreciated.
A locally possessed value, in the sense that matters here, is one \emph{carried by
the particle alone, independently of what happens elsewhere}.
The phrase joins two ideas that philosophers of physics keep apart: that nothing
distant influences the value, which is locality proper, and that the value
belongs to the particle by itself, which is usually called separability.
Such a property belongs to the particle in its own right, at its own location,
whatever apparatus is chosen a light-year away.
It is intrinsic to the particle.

This is exactly the kind of property that classical physics attributed to matter:
the position, momentum, and spin of a particle are properties the particle has,
regardless of what anyone does with it.

There is a moment, for many people who study this argument carefully, when it
really lands.
I remember mine.
I had understood the theorem technically for years before I understood it
philosophically.
The click came when I realized that Bell was not saying the particles are
mysterious or unknowable.
He was not saying that our measurements disturb a reality that was already there.
He was saying that the kind of reality we had assumed --- the kind in which each
particle carries its own definite properties, independent of what happens
elsewhere --- simply does not exist.
Not that it is hidden from us.
Not that it is too small to see.
But that it is not there.

The EPR program is formally closed.
The particles were never carrying a predetermined answer.
The question of what \emph{was} there before measurement is precisely the question
that the rest of this chapter addresses.

\subsection*{What the Theorem Assumes}

Locality and prior definite values are the premises everyone remembers, but the
theorem rests on two more.
One is \emph{measurement independence}: the settings chosen at the two stations
are statistically independent of whatever the particles carry from the source.
A superdeterministic theory denies this, at the price of a correlation, fixed in
the distant past, between the particles and the choices of those who will
measure them; tests that took their settings from the light of distant quasars
push any such correlation back billions of years.
The other is that each measurement has a \emph{single outcome}, which Everett's
many-worlds interpretation denies.
Retrocausal models drop a further, usually implicit assumption: that nothing is
influenced by what happens later.

The premise of prior values does less work than the usual summary suggests.
Bell's original derivation took it from the EPR argument itself: measured along
the same axis, the two particles always give opposite results, so if locality
holds, each result must have been fixed in advance.
In later work Bell recast the theorem to cover theories in which outcomes are
only probable, provided the probabilities at each station depend only on what
lies in that station's past.
Taking the first horn, then, does not by itself keep locality.
Bell concluded that his theorem and the experiments together establish
nonlocality, and a number of philosophers of physics agree with him.
The relational interpretation keeps locality, on its own account, only by a
further step: it denies that the outcome at one station is a fact for the other
until the two records are brought together~\citep{Rovelli2021}.
What the experiments rule out is the world Einstein hoped for, in which separate
things behave according to what each carries and what is near it.


% ─────────────────────────────────────────────────────────────────────────────
\section{Step 2: If No Intrinsic Properties, Then What?}
\label{sec:ch5-step2}
% ─────────────────────────────────────────────────────────────────────────────

Bell's theorem tells us what quantum particles \emph{do not} have: definite,
locally possessed properties prior to measurement.
What, then, do they have?

Here the argument changes character.
Bell's theorem is a negative result, and it does not dictate its own
replacement: Bohm, Everett, and Rovelli all accept the theorem and answer this
question differently.
The answer I defend is relational, and I defend it as the best explanation of
what the experiments show, not as a deduction from them.
It comes from looking carefully at what happens when a measurement \emph{is}
performed.

When a particle interacts with a measuring apparatus, an outcome occurs.
The particle registers as spin-up or spin-down; it lands at one location or another
on the detector screen; it triggers a click in one detector and not the other.
The outcome is definite.
It is real.
After the measurement, there is a fact: \emph{this particle, in this interaction,
with this apparatus, had this result}.

The philosophical weight rests on the italicized structure.
The outcome exists \emph{as a relational fact}: it is a fact about this particle
in this interaction, not a fact about the particle in isolation.
Before the interaction, there was no fact about what the outcome would be ---
not because we did not know, but because the fact did not yet exist.
The measurement event is not the reading of a pre-written text.
It is the event through which the text comes to be written.

This is what Carlo Rovelli calls relational quantum mechanics~\citep{Rovelli1996}.
The central proposal is this: quantum states are always states \emph{relative to}
some system.
There is no absolute, observer-independent quantum state of the universe.
A particle does not simply \emph{have} spin-up.
It has spin-up \emph{relative to this detector, in this interaction, from the
perspective of this measurement event}.

This might sound like a retreat into subjectivism, but it is not.
The outcome is real.
If I measure a particle's spin and get spin-up, and then you compare notes with me,
you will find that I got spin-up.
The definiteness belongs to the relational event, not to the isolated particle.

Two implications of this picture deserve emphasis.

The first is that the measurement event is \emph{constitutive} rather than
\emph{revelatory}.
In classical physics, measurement reveals a pre-existing fact.
The speedometer does not create the car's velocity; it reads what is already
there.
In quantum mechanics, on the relational picture, measurement is different in kind.
It is the event through which the property comes into being.

The second implication is about entanglement.
When two particles interact and become entangled, the correlations between them
are not the result of information passing between them.
They are not the result of a pre-agreed hidden plan.
They are the result of a shared relational history: the two systems' states are
defined partly in terms of each other, because they entered into a relational event
that bound their subsequent behavior together.
Entanglement, on the relational picture, is the persistence of relational
constitution across space and time.

A second theorem bears on this answer more directly than Bell's, and it involves
no distance at all.
In 1967 Simon Kochen and Ernst Specker considered a single particle of spin~1.
The squares of its spin components along any three mutually perpendicular
directions can be measured together, and quantum mechanics requires that, in
suitable units, two of the three results be 1 and the third 0.
Kochen and Specker showed that for a suitably chosen finite set of directions
there is no way to fix each direction's value in advance, independently of its
partners, so that every perpendicular triple comes out right.
If there are values before measurement, the value found for a direction must
depend on what else is measured alongside it, on the measuring context.
This is \emph{contextuality}, and the theorem establishes it for any system with
at least three distinguishable states.
It does not rule out hidden variables; Bohmian mechanics is contextual in just
this sense for every quantity except position.
What it secures is the minimal core of the relational answer: a property of this
kind belongs to the system in a context, not to the system alone.
The fuller claim, that the value comes into being in the interaction, is
interpretation, and Bohm's theory shows that it can be denied.

The wave function, on this reading, is not a hidden description of an absolute
reality.
It is a compact specification of the relational possibilities of a quantum system:
the set of outcomes that can come about when the system enters various kinds of
interactions.
Heisenberg's uncertainty principle~\citep{Heisenberg1927} is not, on this reading,
a limitation imposed on us by our large and clumsy measurement instruments.
It is a statement about the ontological structure of quantum systems: properties
that are jointly definite are constrained, because definiteness is not intrinsic
but relational, and different relational contexts actualize different aspects of
the system's potential.

What is emerging here is a picture of the quantum world that is fundamentally
different from the classical one.
In the classical picture, the world is composed of things with intrinsic
properties.
In the emerging quantum picture, it is composed of relational events, and what we
call a ``particle'' is better understood as a locus of possible interactions, each
of which will actualize some definite aspect of the system's behavior.

There is a philosophical tradition that arrived at a closely similar picture.
It did not arrive at it through quantum mechanics.
It arrived at it through a different route, from first principles of metaphysics,
in the first decades of the twentieth century.
Its name is process philosophy, and its founder was Alfred North Whitehead.
Before I can explain why process ontology is the best-developed framework for
understanding what quantum mechanics has discovered, I need to establish something
more precise: that the framework it replaces, substance ontology in its
classical local form, is not merely inadequate but formally incompatible with the
quantum world.


% ─────────────────────────────────────────────────────────────────────────────
\section{Step 3: Why Local Substance Ontology Is Formally Incompatible}
\label{sec:ch5-step3}
% ─────────────────────────────────────────────────────────────────────────────

The word ``formally'' in this section's title carries the full weight of the
argument.
I am not claiming that substance ontology is philosophically unfashionable,
or that it sits uneasily with modern physics, or that a process-oriented
metaphysician finds it aesthetically unappealing.
I am claiming that there is a valid deductive argument from the experimental
confirmation of Bell inequality violations to the conclusion that substance
ontology, in its classical local form, has no instantiation in the physical world,
given the assumptions set out at the end of Step~1.

Let me state the argument explicitly.

\emph{Premise 1:} Substance ontology requires that fundamental physical entities
possess intrinsic properties --- properties they hold independently of their
relations to other things.
This is not an arbitrary requirement; it is what the word ``substance'' has meant
in the classical tradition.
A substance, in the classical philosophical tradition from Aristotle through
Descartes to Newton, is something that exists in its own right, with its own
definite characteristics.
The billiard ball has mass, charge, and position.
These are properties of the ball.
They belong to it whether or not anyone measures it, whether or not it is
interacting with anything else.
A substance that had no properties in isolation would not be a substance in any
recognizable sense; it would be an empty grammatical subject without predicates.
For the argument, Premise~1 has to be read the way classical physics read it:
the intrinsic properties include those that fix how a thing responds to
measurement, such as its position and its spin along an axis, and each thing's
properties are its own, so that the state of a pair is settled by the states of
its members.
Not every property is in question.
An electron's mass and charge are the same in every state it can occupy, and
nothing in Bell's theorem touches them; the argument concerns the properties that
vary from state to state.

\emph{Premise 2:} Bell's theorem, confirmed in experiments since the early 1970s,
proves that no local theory can attribute definite properties to quantum particles
prior to measurement.
More precisely: no theory in which each particle's outcome is fixed, or given its
probability, by locally possessed hidden variables can reproduce the statistical
correlations that quantum mechanics predicts and experiment confirms, so long as
the settings are independent of those variables and each measurement has a
single outcome.

\emph{Conclusion:} Therefore, the fundamental physical entities of quantum
mechanics cannot be substances in the classical sense.
They do not possess intrinsic, locally held values for the properties that
measurements reveal.
Substance ontology --- understood as the claim that reality is fundamentally
composed of separate entities with intrinsic local properties --- is formally
ruled out by quantum mechanics.

The logic is valid.
If you accept both premises, you must accept the conclusion.
Neither premise is beyond dispute, though.
Premise~2 carries its auxiliary assumptions with it, and each has been denied by
someone serious.
Premise~1 describes the classical picture, not every philosophy that has used
the word ``substance.''
A neo-Aristotelian, for whom relation is one of the categories, could treat an
electron's spin along an axis as a relational accident of a substance whose
charge is essential to it.
That reply is fair, and it narrows the target to the picture Einstein defended
against quantum mechanics: a world in which each thing carries its whole physical
state itself, changed only by what is near it.

Let me address the most serious ways of resisting this conclusion.

\subsection*{The Bohmian Objection}

The most direct challenge to the conclusion comes from Bohmian mechanics, also
called the pilot-wave interpretation~\citep{Bohm1980}.
In Bohm's framework, quantum particles do have definite positions at all times.
The wave function does not collapse; it guides the particles along definite
trajectories through a globally defined pilot wave.
Measurement outcomes are definite because the particles were at definite
locations throughout; the uncertainty of quantum mechanics reflects our ignorance
of the precise initial conditions, not a genuine ontological indefiniteness.

If Bohmian mechanics is right, has substance ontology been saved?

In part, and this objection sets the scope of the whole chapter.

Bell's theorem eliminates local hidden variables, not hidden variables as such.
Bohmian mechanics survives Bell precisely by abandoning locality: the pilot wave
is defined on the full configuration space of the system, not on three-dimensional
physical space.
The velocity of each particle depends at every instant on the positions of all the
particles with which it is entangled, however far away they are, so that setting
the apparatus at A one way rather than another changes the velocity of particle B
at once.
No one can use this to signal, but within the theory the influence is real.
Relativistic versions of the theory need a preferred slicing of spacetime into
moments of simultaneity, a structure no experiment can detect and relativity was
built to do without.

What the Bohmian keeps is a fact, at every moment, about where each particle is,
a fact that concerns the particle alone.
That is recognizably a kind of substance.
It is tempting to reply that a particle whose path depends on everything else has
no intrinsic properties, but the reply proves too much: in Newton's theory a
planet's path depends at every instant on every other mass in the universe, and
no one concluded that Newton's planets were relational events.
What the Bohmian gives up is the rest of the classical picture.
Spin and momentum are not quantities the particle carries; what a measurement of
them finds depends on the arrangement of the whole apparatus, so that the same
particle in the same state can give different answers in different experimental
contexts.
And the law that moves the particles is written in a wave function that belongs
to no particle and to no region of space, and that for entangled systems cannot
be divided into states of the parts.

The fair verdict is that quantum mechanics does not rule out every substance
ontology.
It rules out the local one, and whatever substance survives has position as its
only intrinsic property and is moved by a nonlocal law that picks out a
simultaneity relativity does not recognize.
A reader willing to pay that price has a coherent alternative to the view of this
book.
My reasons for not paying it are comparative: Bohm privileges position over all
the other quantities the formalism treats alike, adds to the wave function a
second kind of entity whose work is to make outcomes single, and still needs a
holistic structure, the wave function, of the kind process ontology takes as
basic.
The Bohmian objection leaves the argument's first half standing and marks its
boundary.

\subsection*{Many Worlds, QBism, and the Whole}

Two further rivals escape Bell's theorem by denying assumptions other than
locality.
The Everett interpretation, usually called many-worlds, keeps the wave function
and nothing else.
It never collapses; a measurement entangles apparatus with system, and every
outcome occurs, each in its own branch of one universal quantum state, so that no
influence need pass between the stations.
The price is familiar: branches that multiply with every measurement-like
interaction, and a long-standing difficulty in saying what the Born-rule
probabilities are probabilities of when every outcome occurs.
Everett does save one substance.
The universal state is absolute, and on this view it is the whole of physical
reality, but its parts have no states independent of one another, only states
relative to the rest, which is why Everett called his proposal the relative-state
formulation.
Some metaphysicians, Jonathan Schaffer among them, have taken entanglement to
support monism, on which the cosmos as a whole is the one fundamental thing.
That is a coherent substance ontology, but it keeps an intrinsic nature for the
whole alone and makes each part what it is through its relations within it,
which I count as a concession to the relational view.

QBism, developed by Christopher Fuchs and his collaborators, reads the quantum
state as an agent's degrees of belief about the consequences of her actions, and
a measurement outcome as that agent's experience.
It escapes Bell's theorem by declining to treat the outcomes at the two stations
as facts for everyone at once: a correlation becomes a fact for an agent when she
has both results in hand.
It restores no substance; what it affirms, under the name of participatory
realism, is that each measurement is a creative encounter between agent and
world, which a process philosopher will find congenial.
Its price mirrors Everett's: Everett offers a world without single outcomes,
QBism outcomes without a third-person description of the world they occur in.

Each rival, then, denies a different premise.
None keeps the classical picture of separate things, each carrying its whole
state, and that common result is all the deductive half of this chapter claims.
Which rival to prefer is a further question, and its answer is not a theorem.

\subsection*{The Macroscopic-Scale Intuition}

A different resistance comes not from physics but from everyday experience.
The table I am writing on has definite mass, a definite location, a definite
surface area.
These seem like paradigm cases of intrinsic properties.
The mass is not in question, for the reason given under Premise~1, but the
table's location and the settled arrangement of its parts are exactly the kind of
property Bell's theorem concerns.
If substance ontology is formally ruled out, how do we explain the table?

The answer comes from decoherence~\citep{Zurek2003}.

In quantum mechanics, a system becomes classical-appearing when it interacts
strongly with its environment.
The environment --- air molecules, photons, gravitational fluctuations ---
continuously entangles with the system, and this entanglement very rapidly
suppresses the system's local quantum coherence (its ability to show
interference between distinct classical states), the feature that distinguishes
quantum behavior from classical behavior.
Strictly, the coherence is not destroyed but exported into correlations between
the system and its surroundings, as the fourth alignment below explains.
For a macroscopic object like a table, the decoherence process happens so fast
and so thoroughly that no superposition of distinct locations survives long
enough to be observable.
The table \emph{appears} to have a definite location because its environment
keeps recording where it is, which confines it to a small set of robust states
between which no interference can be seen.
Decoherence explains why those alternatives stop interfering; it does not by
itself explain why one of them occurs and not another.

But decoherence does not restore substance ontology.
The table's apparent definiteness is a consequence of its interactions, not
evidence that it has intrinsic properties independent of interaction.

Our intuition that reality is made of things with definite properties is not
wrong.
It is a very accurate representation of how reality behaves at the scales we
evolved to perceive.
But it is not a description of the fundamental structure of reality.
It is what a relational quantum world looks like when viewed through the
decoherence filter of macroscopic observation.
At the quantum scale, in its local form, it fails --- and Bell's theorem tells us
why.


% ─────────────────────────────────────────────────────────────────────────────
\section{Step 4: Process Ontology as the Best-Fitting Framework}
\label{sec:ch5-step4}
% ─────────────────────────────────────────────────────────────────────────────

With local substance ontology formally ruled out, the question becomes: what
framework takes its place?

Process ontology, most fully developed by Alfred North Whitehead in
\textit{Process and Reality}~\citep{Whitehead1929} and extended by subsequent
philosophers of physics including \citet{Epperson2004} and
\citet{Stapp2007}, provides the most systematic answer that has so far
been developed.

The central inversion of process ontology is this: not things, but events.
Not persistent objects with properties, but relational events through which
properties come to be.
What we call ``things'' are, in Whitehead's framework, patterns of recurring
relational events --- high-frequency societies of ``actual occasions'' that
maintain a recognizable structure over time.
The stone, the electron, the organism, the mind: each is a pattern of relational
becoming, not an independent substance.

To show why I take this framework to be not merely an alternative to substance
ontology, but the most accurate philosophical description of what quantum
mechanics has discovered, I want to map the four features of quantum mechanics
that have proven most resistant to classical interpretation onto the four central
categories of Whitehead's framework.

\subsection*{First Alignment: Property Indefiniteness and Actual Occasions}

The hardest feature of quantum mechanics to swallow for anyone trained in classical
physics is property indefiniteness: the fact that a quantum particle does not have
a definite spin, position, or momentum before it is measured.
This is not a limitation of our instruments.
It is, on the reading I have defended, a feature of reality; what Bell proved is
that it cannot be explained away by local hidden values.

Substance ontology has no place for this.
If reality is composed of substances with intrinsic properties, every substance
has its properties at every moment.
The question ``what is the particle's spin before measurement?'' must have an
answer.
Bell showed that it has no local answer, and Kochen and Specker that it has no
answer independent of the measuring context.

Process ontology, by contrast, has exactly the right structure to accommodate
property indefiniteness.
An actual occasion, in Whitehead's account, begins as pure potentiality.
It is in the process of \emph{becoming} --- of constituting itself through its
relational grasping of what came before.
Definiteness is not given; it is achieved.
The occasion achieves its definite character through \emph{concrescence}, the
process by which potential is converted into actuality through relational
engagement.

The quantum particle before measurement maps onto the actual occasion in its
pre-concrescence phase.
It is not indefinite because of our ignorance.
It is indefinite because definiteness is not the ontological ground state of
relational events --- it is the outcome of a relational process.
The wave function, on this reading, is a precise description of what a
pre-concrescence quantum occasion actually is: a structured distribution of
relational potentials, waiting to be actualized through interaction.

This identification was developed in formal detail by \citet{Epperson2004} and
constitutes one of the strongest specific alignments between the two frameworks.
Epperson means it as more than an analogy: as a claim that the structure of
actual occasions --- potential, process, actualization --- is the correct
description of the ontological structure that quantum mechanics has
experimentally revealed.

\subsection*{Second Alignment: Entanglement and Universal Prehension}

Entanglement is the most distinctive feature of quantum mechanics and the one
that most decisively breaks the classical picture.
Two particles, once they have interacted, share a quantum state that cannot be
written as a simple product of their individual states.
Measurements on the two are correlated, however far apart they are, in a way
that neither experimenter can use to send a signal and that, Bell proved, no
prior agreement or local hidden variable can explain.
The correlations are a feature of the relational structure of the quantum world
itself.

Substance ontology has no adequate account of entanglement.
If each particle is a substance with its own intrinsic properties, its state
should be fully specifiable independently of any other particle.
Entanglement says otherwise: the state of each particle is, in part, defined by
its relational history with the other.
The particles are not separable substances; they are aspects of a single
relational structure.

Process ontology built this feature into its foundations.
The central act in Whitehead's ontology is \emph{prehension}: the relational
grasping by which each actual occasion incorporates prior occasions into its own
becoming.
Prehension is not perception in the human sense.
It is the most fundamental form of relational responsiveness in the universe,
operative at every scale.
And it is \emph{universal}: every actual occasion prehends, to some degree, every
prior occasion.
There is no quantum event that is fully isolated from its relational history.

Entanglement, on the process-ontological reading, is the quantum formalism's
expression of what universal prehension looks like at the subatomic scale.
Two particles that once interacted have entered into a prehensive relationship;
their subsequent states are defined partly by that relationship.
The correlations that Bell proved cannot be explained by local hidden variables
are, on this reading, the expected consequence of a world in which relations
are ontologically prior to relata --- in which things do not first exist and then
relate, but rather come into being through their relational constitution.

One version of this reading, though, is refuted by Bell's theorem itself.
Suppose each particle carried forward its own prehension of the shared past, and
each outcome were settled by that inheritance together with an act of
self-determination at its own station.
The shared past would then be exactly the kind of common cause Bell's theorem
excludes, and a world built that way would obey Bell's inequality.
The relational reading works only if the entangled pair is one relational whole,
whose actualization at one station is not independent of what happens at the
other.
What that asks of Whitehead's scheme I take up at the end of this section.

\subsection*{Third Alignment: Observer-Dependence and Perspectival Prehension}

Quantum states are, on Rovelli's relational interpretation, always relative to
some observer or system.
Two observers who have not compared notes may, with perfect theoretical legitimacy,
assign different states to the same system, because quantum states are facts
about relations, and different observers stand in different relational contexts.

This feature of quantum mechanics has seemed deeply uncomfortable to many
physicists, because it appears to threaten objectivity.
If the state is always relative to an observer, is there any objective reality?

Process ontology dissolves the discomfort without sacrificing objectivity.
In Whitehead's framework, every actual occasion is irreducibly \emph{perspectival}.
It prehends the world from its own particular location in the relational network,
incorporating those prior occasions that are relevant to its own becoming.
Every prehension is a view from somewhere.

And yet the results are objective.
Once an occasion has achieved its definite character through concrescence, it
becomes part of the objective past available for prehension by all subsequent
occasions.
The definiteness of the past is absolute; only the perspectives of observers in
the present are multiple.
Rovelli's relational quantum mechanics shares half of this structure: its facts
about measurement outcomes are as real as any facts can be, and the quantum
states that precede them are perspectival.
It does not share the other half.
For Rovelli a fact relative to one system need not yet be a fact relative to
another, as the case of Wigner's friend in Chapter~\ref{ch:observer-dependence}
will show, whereas Whitehead's perished occasions are the same for every occasion
that inherits them.

Observer-dependence in quantum mechanics is not a feature requiring apology.
It is the quantum signature of what a perspectival, relational universe looks like
when described mathematically.

\subsection*{Fourth Alignment: Decoherence and the Objective Past}

The transition from quantum to classical behavior is described by the theory of
decoherence~\citep{Zurek2003}.
When a quantum system interacts with a large environment, the correlations
between system and environment spread into the environment in a way that is
practically irreversible, and to any local observer the system appears to have
acquired a definite classical state.
The quantum coherence has not been destroyed; it has been distributed throughout
the environment beyond any practical possibility of recovery.

The measurement problem, how a definite outcome emerges from a superposition, is
a sharper difficulty than the transition itself, and it is not peculiar to
substance ontology.
It arises for any view that takes the quantum state to be complete, lets it
always evolve smoothly by the Schr\"odinger equation, and expects single
outcomes.
Bohm gives up the first of these, collapse theories the second, and Everett the
third; process readings such as Stapp's side with the collapse theories, treating
actualization as a real event.

Process ontology has a natural account of the transition.
In Whitehead's framework, each actual occasion undergoes concrescence --- the
process of achieving its definite character through relational engagement --- and
then \emph{perishes}.
It passes from the subject of experience to the superject: the completed fact
that is now available for prehension by subsequent occasions.
The perished occasion enters the objective past.
It is no longer in process; it is finished, definite, fixed.

Decoherence, on the process-ontological reading, is what the perishing of actual
occasions looks like from the perspective of quantum mechanics.
As a quantum system interacts with its environment, the complexity of its
relational entanglement grows until the system effectively completes its
concrescence: it becomes a finished fact, a fixed element of the relational past,
no longer exhibiting the superposition behavior characteristic of pre-concrescence
occasions.
On this reading the transition from quantum to classical is the transition from
the potentiality phase to the fixed objective past.

The fit is not exact, and the gap is instructive.
Decoherence is gradual and in principle reversible, and it leaves a mixture of
all the robust alternatives, not one of them; concrescence is a completed act
that yields one.
Decoherence supplies the menu, the alternatives that can become definite and the
reason they no longer interfere, and the process reading adds the choice.
That addition goes beyond the physics, as Bohm's positions and the collapse
postulate also do.
\citet{Epperson2004} develops this mapping in technical detail, arguing that
Whitehead's categories of potential phase, concrescence, and superjecthood
correspond naturally to the quantum-mechanical description of pre-measurement,
measurement, and post-measurement phases.
He presents the alignment as structural, not approximate, and within the limit
just described I think he is right.

\subsection*{The Overall Picture}

The four alignments are not independent.
They all flow from the same fundamental inversion: process ontology puts events
and relations where substance ontology put things and properties.
This single inversion explains property indefiniteness (definiteness is the
outcome of relational process, not a given), entanglement (relations are
ontologically prior to relata), observer-dependence (every prehension is
perspectival), and the quantum-to-classical transition (occasions perish into
a fixed objective past).

\citet{Stapp2007} approaches the same convergence from the direction of quantum
mechanics itself, arguing that the structure of quantum events --- with their
abrupt transitions from potential to actual --- is naturally described in the
ontological framework that Whitehead developed.
\citet{Griffin2000} extends the process framework to consciousness and
experience.

I do not think the convergence is accidental, and it is not forced.
Process ontology was not developed by reading quantum mechanics and working
backward to a metaphysics.
It was developed from first principles by a philosopher-mathematician who saw
the inadequacy of substance ontology well before quantum mechanics had made
that inadequacy experimentally precise.
That independence needs one qualification.
Whitehead knew the physics of his day: he had written on relativity, and he
discussed the older quantum theory in \textit{Science and the Modern World}.
But \textit{Process and Reality} makes little direct use of the new quantum
mechanics of 1925--27, and it appeared thirty-five years before Bell's theorem.
That Whitehead's framework fits the post-Bell picture with this degree of
structural precision is itself philosophically significant.
It suggests that both the physics and the philosophy were detecting the same
underlying feature of reality, by different methods.

\subsection*{What Process Ontology Must Still Pay}

Process ontology has costs of its own, and fairness requires listing them.

The first was flagged under the second alignment.
In Whitehead's scheme an occasion inherits only from its past, and contemporary
occasions, those outside each other's causal past, are causally independent.
Read literally, that makes his cosmology local in Bell's sense, so the scheme has
to be amended where this chapter has leaned on it hardest.
Stapp's Whiteheadian reading lets each actual event reshape the state of the
whole system at once and orders actual events in one global sequence, close kin
to the preferred slicing of spacetime that Bohm needs.
The alternative follows Rovelli: keep locality, and deny that the outcome at one
station is a fact for the other before the records meet, which gives up the
feature noted under the third alignment, a past that is the same for every
occasion that inherits it.
Process ontology must take one route or the other, and neither is free.

The second is the choice conceded under the fourth alignment: concrescence
supplies the single outcome that decoherence does not, and that is an addition to
the physics.
I think it better motivated than its rivals, since the metaphysics made it for
its own reasons, but it is an addition all the same.

The third is experience.
Whitehead's occasions have an experiential structure all the way down, in the
generalized sense described in Chapter~\ref{ch:process-ontology}, and nothing in
quantum mechanics requires that.
A reader can accept every alignment in this section while declining it.

With those costs on the table, I still judge process ontology the best fit: it
takes as basic the relational holism each rival needs anyway, it treats all
observables alike, and it gives single outcomes a place without multiplying
worlds.
A careful reader may weigh the costs differently.
I will return to the question of independent detection of the same structure at
length in the next chapter.
For now, I want to collect the argument of this chapter into a precise conclusion.


% ─────────────────────────────────────────────────────────────────────────────
\section{The Strong Conclusion}
\label{sec:ch5-conclusion}
% ─────────────────────────────────────────────────────────────────────────────

Let me be as precise as possible about what has been established and what has not.

\textbf{What has been established:}

\emph{First}, the experimental record of Bell inequality violations, culminating in
the loophole-free tests of 2015--2017 and recognized by the 2022 Nobel Prize,
confirms that no local hidden-variable theory can reproduce the predictions of
quantum mechanics, provided the settings are independent of the hidden variables
and each measurement has a single outcome.
This is not a philosophical preference or an interpretation.
It is an experimentally confirmed mathematical theorem.

\emph{Second}, local substance ontology requires exactly what Bell eliminates:
intrinsic, locally possessed values for the properties that measurements reveal.
Therefore local substance ontology is formally incompatible with quantum mechanics.
This conclusion follows by deduction from the commitments of the classical
picture and the experimental confirmation of Bell's theorem, and it is the part
of the argument that earns the word ``requires'' without qualification.

\emph{Third}, each surviving alternative to local substance ontology gives
relations, or the whole, a role the classical picture denied them: Bohm's
nonlocal wave function belongs to no particle, the parts of Everett's universal
state have only relative states, and QBism locates each outcome in an agent's
encounter with the world.
The relational interpretation makes the point explicit: quantum states are facts
about relations between systems.
Process ontology provides the most fully developed metaphysical framework for this
relational picture.

\emph{Fourth}, four central features of quantum mechanics --- property indefiniteness,
entanglement, observer-dependence, and the quantum-to-classical transition via
decoherence --- align structurally with four features of process
ontology: the potentiality phase of actual occasions, universal prehension,
perspectival concrescence, and the perishability of occasions into the objective
past.
Epperson and Stapp develop these alignments in technical detail, as more than
analogies.
They are contested, and some hold only with the amendments described at the end
of Step~4.

\textbf{What has not been established:}

I have not proven that process ontology is true.
No argument of this kind can do that.
What has been established is an asymmetry: one framework has been formally ruled
out, and the remaining frameworks each give relations a role that process
ontology most fully develops.
This asymmetry is epistemically significant even in the absence of proof.

Nor have I refuted the rivals.
Bohmian mechanics, the Everett interpretation, and QBism remain coherent; that
each pays more than process ontology for what it buys is a comparative judgment.
The deductive half of the argument also inherits Bell's assumptions, and a reader
who accepts superdeterminism or influences running backward in time can escape
it, though few have found that price worth paying.

I have also not claimed that Whitehead had the complete picture.
Process ontology as Whitehead developed it was not designed with quantum mechanics
in mind; its doctrine that contemporaries are causally independent must be
amended before it can accommodate entanglement, and there are technical aspects
of quantum field theory and the measurement problem that require further
philosophical work.
What I am claiming is structural: the architecture of process ontology maps onto
the architecture of quantum reality with a precision that, as far as I know, no
other general metaphysical system matches.

Finally, I have not resolved the full measurement problem --- the precise account
of how a single definite outcome emerges from a superposition in any given
measurement event.
That remains one of the deepest open questions in the foundations of physics.
What process ontology provides is not a solution to the measurement problem but
a metaphysical framework in which the problem takes a natural form.
The question ``how does potential become actual through relational engagement?''
is not an embarrassment within process ontology.
It is precisely the question that the framework was built to ask.

The conclusion of this chapter, stated as precisely as I can manage: local
substance ontology is formally ruled out by quantum mechanics, given the
assumptions of Bell's theorem.
Process ontology is the most fully developed alternative, and four structural
features of the quantum world align closely with its central categories.
Quantum mechanics therefore requires something like process ontology, an
ontology in which relations do work the intrinsic properties of separate things
cannot, and among the ontologies of that kind, process ontology is the best fit.
That is a powerful asymmetry, and it has a direct implication for the
conversations that occupy the rest of this book.

There is, it turns out, a third tradition that arrived at the same relational
picture of reality.
It did not arrive through mathematics, and it did not arrive through experiment.
It arrived through a method of its own, in writings composed between the 1860s and
1921 --- before Whitehead named the framework and more than four decades before
Bell proved the theorem that ruled out its local rival.
That is the subject of the next chapter.

\begin{convergencemoment}{Why Quantum Mechanics Points to Process Ontology}
The experimental record of Bell inequality violations, confirmed from Aspect's
1982 tests through the 2022 Nobel Prize, formally rules out local substance
ontology: the classical picture of particles with intrinsic, locally possessed
properties has no instantiation in the quantum world.

Process philosophy, developed by Whitehead in 1929, thirty-five years before
Bell's theorem, offers the most fully developed alternative: a framework in which
reality consists not of things with properties but of relational events through
which properties come to be.
The four central features of quantum mechanics that resist classical explanation
--- property indefiniteness, entanglement, observer-dependence, and the
quantum-to-classical transition --- map onto four features of
process ontology with a structural precision that is hard to credit to
coincidence.

From the \bahai side: three texts composed between the 1860s and 1921 articulate
a process ontological picture of reality --- existence as relational affinity,
creation as relational event, process as the fundamental mode of being ---
independently of and prior to both Whitehead and quantum mechanics.
Three independent instruments; one structure.
\end{convergencemoment}
